\section{Relative normalization}
\label{section-normalization-X-in-Y}

\noindent
In this section we construct the {\it normalization of one scheme in another}.

\begin{lemma}
\label{lemma-integral-closure}
Let $X$ be a scheme. Let $\mathcal{A}$ be a quasi-coherent sheaf
of $\mathcal{O}_X$-algebras. The subsheaf $\mathcal{A}' \subset \mathcal{A}$
defined by the rule
$$
U \longmapsto \{f \in \mathcal{A}(U) \mid
f_x \in \mathcal{A}_x \text{ integral over } \mathcal{O}_{X, x}
\text{ for all }x \in U\}
$$
is a quasi-coherent $\mathcal{O}_X$-algebra, the stalk $\mathcal{A}'_x$
is the integral closure of $\mathcal{O}_{X, x}$ in $\mathcal{A}_x$, and
for any affine open $U \subset X$ the ring
$\mathcal{A}'(U) \subset \mathcal{A}(U)$ is
the integral closure of $\mathcal{O}_X(U)$ in $\mathcal{A}(U)$.
\end{lemma}

\begin{proof}
This is a subsheaf by the local nature of the conditions.
It is an $\mathcal{O}_X$-algebra by
Algebra, Lemma \ref{algebra-lemma-integral-closure-is-ring}.
Let $U \subset X$ be an affine open. Say $U = \Spec(R)$
and say $\mathcal{A}$ is the quasi-coherent sheaf associated to
the $R$-algebra $A$. Then according to
Algebra, Lemma \ref{algebra-lemma-integral-closure-stalks}
the value of $\mathcal{A}'$ over $U$ is given by the integral
closure $A'$ of $R$ in $A$. This proves the last assertion of
the lemma. To prove that $\mathcal{A}'$ is quasi-coherent, it
suffices to show that $\mathcal{A}'(D(f)) = A'_f$. This follows
from the fact that integral closure and localization commute, see
Algebra, Lemma \ref{algebra-lemma-integral-closure-localize}.
The same fact shows that the stalks are as advertised.
\end{proof}

\begin{definition}
\label{definition-integral-closure}
Let $X$ be a scheme. Let $\mathcal{A}$ be a quasi-coherent sheaf
of $\mathcal{O}_X$-algebras. The {\it integral closure of $\mathcal{O}_X$
in $\mathcal{A}$} is the quasi-coherent $\mathcal{O}_X$-subalgebra
$\mathcal{A}' \subset \mathcal{A}$ constructed in
Lemma \ref{lemma-integral-closure} above.
\end{definition}

\noindent
In the setting of the definition above we can consider the morphism
of relative spectra
$$
\xymatrix{
Y = \underline{\Spec}_X(\mathcal{A}) \ar[rr] \ar[rd] & &
X' = \underline{\Spec}_X(\mathcal{A}') \ar[ld] \\
& X &
}
$$
see Lemma \ref{lemma-affine-equivalence-algebras}.
The scheme $X' \to X$ will be the normalization of $X$ in the scheme $Y$.
Here is a slightly more general setting. Suppose we have a
quasi-compact and quasi-separated morphism $f : Y \to X$
of schemes. In this case the sheaf of
$\mathcal{O}_X$-algebras $f_*\mathcal{O}_Y$ is quasi-coherent, see
Schemes, Lemma \ref{schemes-lemma-push-forward-quasi-coherent}.
Taking the integral closure $\mathcal{O}' \subset f_*\mathcal{O}_Y$
we obtain a quasi-coherent sheaf of $\mathcal{O}_X$-algebras
whose relative spectrum is the normalization of $X$ in $Y$. Here is
the formal definition.

\begin{definition}
\label{definition-normalization-X-in-Y}
Let $f : Y \to X$ be a quasi-compact and quasi-separated morphism of schemes.
Let $\mathcal{O}'$ be the integral closure of $\mathcal{O}_X$ in
$f_*\mathcal{O}_Y$. The {\it normalization of $X$ in $Y$} is the
scheme\footnote{The scheme $X'$ need not be normal, for example if
$Y = X$ and $f = \text{id}_X$, then $X' = X$.}
$$
\nu : X' = \underline{\Spec}_X(\mathcal{O}') \to X
$$
over $X$. It comes equipped with a natural factorization
$$
Y \xrightarrow{f'} X' \xrightarrow{\nu} X
$$
of the initial morphism $f$.
\end{definition}

\noindent
The factorization is the composition of the canonical morphism
$Y \to \underline{\Spec}_X(f_*\mathcal{O}_Y)$ (see
Constructions, Lemma
\ref{constructions-lemma-canonical-morphism})
and the morphism of relative spectra coming from the inclusion map
$\mathcal{O}' \to f_*\mathcal{O}_Y$. We can characterize the
normalization as follows.

\begin{lemma}
\label{lemma-characterize-normalization}
Let $f : Y \to X$ be a quasi-compact and quasi-separated morphism of schemes.
The factorization $f = \nu \circ f'$, where $\nu : X' \to X$ is the
normalization of $X$ in $Y$ is characterized by the following
two properties:
\begin{enumerate}
\item the morphism $\nu$ is integral, and
\item for any factorization $f = \pi \circ g$, with $\pi : Z \to X$
integral, there exists a commutative diagram
$$
\xymatrix{
Y \ar[d]_{f'} \ar[r]_g & Z \ar[d]^\pi \\
X' \ar[ru]^h \ar[r]^\nu & X
}
$$
for some unique morphism $h : X' \to Z$.
\end{enumerate}
Moreover, the morphism $f' : Y \to X'$ is dominant and in (2) the
morphism $h : X' \to Z$ is the normalization of $Z$ in $Y$.
\end{lemma}

\begin{proof}
Let $\mathcal{O}' \subset f_*\mathcal{O}_Y$ be the integral closure of
$\mathcal{O}_X$ as in Definition \ref{definition-normalization-X-in-Y}.
The morphism $\nu$ is integral by construction, which proves (1).
Assume given a factorization $f = \pi \circ g$ with $\pi : Z \to X$
integral as in (2). By Definition \ref{definition-integral}
$\pi$ is affine, and hence $Z$ is the relative
spectrum of a quasi-coherent sheaf of $\mathcal{O}_X$-algebras $\mathcal{B}$.
The morphism $g : Y \to Z$ corresponds to a map of $\mathcal{O}_X$-algebras
$\chi : \mathcal{B} \to f_*\mathcal{O}_Y$. Since $\mathcal{B}(U)$ is
integral over $\mathcal{O}_X(U)$ for every affine open $U \subset X$
(by Definition \ref{definition-integral})
we see from Lemma \ref{lemma-integral-closure}
that $\chi(\mathcal{B}) \subset \mathcal{O}'$.
By the functoriality of the relative spectrum
Lemma \ref{lemma-affine-equivalence-algebras}
this provides us with a unique morphism
$h : X' \to Z$. We omit the verification that the diagram commutes.

\medskip\noindent
It is clear that (1) and (2) characterize the factorization $f = \nu \circ f'$
since it characterizes it as an initial object in a category.

\medskip\noindent
From the universal property in (2) we see that $f'$ does not factor through
a proper closed subscheme of $X'$. Hence the scheme theoretic image of $f'$
is $X'$. Since $f'$ is quasi-compact (by Schemes, Lemma
\ref{schemes-lemma-quasi-compact-permanence} and the fact that $\nu$ is
separated as an affine morphism) we see that $f'(Y)$ is dense in $X'$.
Hence $f'$ is dominant.

\medskip\noindent
Observe that $g$ is quasi-compact and quasi-separated by
Schemes, Lemmas \ref{schemes-lemma-compose-after-separated} and
\ref{schemes-lemma-quasi-compact-permanence}. Thus the last statement
of the lemma makes sense.
The morphism $h$ in (2) is integral by Lemma \ref{lemma-finite-permanence}.
Given a factorization $g = \pi' \circ g'$ with $\pi' : Z' \to Z$
integral, we get a factorization $f = (\pi \circ \pi') \circ g'$ and
we get a morphism $h' : X' \to Z'$. Uniqueness implies that
$\pi' \circ h' = h$. Hence the characterization (1), (2) applies
to the morphism $h : X' \to Z$ which gives the final assertion of the lemma.
\end{proof}

\begin{lemma}
\label{lemma-functoriality-normalization}
Let
$$
\xymatrix{
Y_2 \ar[d]_{f_2} \ar[r] & Y_1 \ar[d]^{f_1} \\
X_2 \ar[r] & X_1
}
$$
be a commutative diagram of morphisms of schemes.
Assume $f_1$, $f_2$ quasi-compact and quasi-separated.
Let $f_i = \nu_i \circ f_i'$, $i = 1, 2$
be the canonical factorizations, where $\nu_i : X_i' \to X_i$ is
the normalization of $X_i$ in $Y_i$. Then there exists a unique
arrow $X'_2 \to X'_1$ fitting into a
commutative diagram
$$
\xymatrix{
Y_2 \ar[d]_{f_2'} \ar[r] & Y_1 \ar[d]^{f_1'} \\
X_2' \ar[d]_{\nu_2} \ar[r] & X_1' \ar[d]^{\nu_1} \\
X_2 \ar[r] & X_1
}
$$
\end{lemma}

\begin{proof}
By Lemmas \ref{lemma-characterize-normalization} (1)
and \ref{lemma-base-change-finite}
the base change $X_2 \times_{X_1} X'_1 \to X_2$
is integral. Note that $f_2$ factors through this morphism.
Hence we get a unique morphism
$X'_2 \to X_2 \times_{X_1} X'_1$ from
Lemma \ref{lemma-characterize-normalization} (2).
This gives the arrow $X'_2 \to X'_1$ fitting into
the commutative diagram and uniqueness follows as well.
\end{proof}

\begin{lemma}
\label{lemma-normalization-localization}
Let $f : Y \to X$ be a quasi-compact and quasi-separated morphism of schemes.
Let $U \subset X$ be an open subscheme and set $V = f^{-1}(U)$.
Then the normalization of $U$ in $V$ is the inverse image of $U$
in the normalization of $X$ in $Y$.
\end{lemma}

\begin{proof}
Clear from the construction.
\end{proof}

\begin{lemma}
\label{lemma-normalization-is-normalization}
Let $f : Y \to X$ be a quasi-compact and quasi-separated morphism of schemes.
Let $X'$ be the normalization of $X$ in $Y$. Then the normalization of
$X'$ in $Y$ is $X'$.
\end{lemma}

\begin{proof}
If $Y \to X'' \to X'$ is the normalization of $X'$ in $Y$, then
we can apply Lemma \ref{lemma-characterize-normalization}
to the composition $X'' \to X$ to get a canonical morphism
$h : X' \to X''$ over $X$. We omit the verification that the
morphisms $h$ and $X'' \to X'$ are mutually inverse (using uniqueness
of the factorization in the lemma).
\end{proof}

\begin{lemma}
\label{lemma-normalization-in-reduced}
Let $f : Y \to X$ be a quasi-compact and quasi-separated morphism of schemes.
Let $X' \to X$ be the normalization of $X$ in $Y$. If $Y$ is reduced, so
is $X'$.
\end{lemma}

\begin{proof}
This follows from the fact that a subring of a reduced ring is reduced.
Some details omitted.
\end{proof}

\begin{lemma}
\label{lemma-normalization-generic}
Let $f : Y \to X$ be a quasi-compact and quasi-separated morphism of schemes.
Let $X' \to X$ be the normalization of $X$ in $Y$. Every generic point of
an irreducible component of $X'$ is the image of a generic point of
an irreducible component of $Y$.
\end{lemma}

\begin{proof}
By Lemma \ref{lemma-normalization-localization} we may assume $X = \Spec(A)$
is affine. Choose a finite affine open covering $Y = \bigcup \Spec(B_i)$.
Then $X' = \Spec(A')$ and the morphisms $\Spec(B_i) \to Y \to X'$
jointly define an injective $A$-algebra map $A' \to \prod B_i$.
Thus the lemma follows from
Algebra, Lemma \ref{algebra-lemma-injective-minimal-primes-in-image}.
\end{proof}

\begin{lemma}
\label{lemma-normalization-in-disjoint-union}
Let $f : Y \to X$ be a quasi-compact and quasi-separated morphism of schemes.
Suppose that $Y = Y_1 \amalg Y_2$ is a disjoint union of two schemes.
Write $f_i = f|_{Y_i}$. Let $X_i'$ be the normalization of $X$ in $Y_i$.
Then $X_1' \amalg X_2'$ is the normalization of $X$ in $Y$.
\end{lemma}

\begin{proof}
In terms of integral closures this corresponds to the following fact:
Let $A \to B$ be a ring map. Suppose that $B = B_1 \times B_2$.
Let $A_i'$ be the integral closure of $A$ in $B_i$. Then
$A_1' \times A_2'$ is the integral closure of $A$ in $B$.
The reason this works is that the elements $(1, 0)$ and $(0, 1)$ of $B$
are idempotents and hence integral over $A$. Thus the integral closure
$A'$ of $A$ in $B$ is a product and it is not hard to see that the factors
are the integral closures $A'_i$ as described above (some details
omitted).
\end{proof}

\begin{lemma}
\label{lemma-normalization-in-universally-closed}
Let $f : X \to S$ be a quasi-compact, quasi-separated and
universally closed morphisms of schemes.
Then $f_*\mathcal{O}_X$ is integral over $\mathcal{O}_S$. In other
words, the normalization of $S$ in $X$ is equal to the factorization
$$
X \longrightarrow \underline{\Spec}_S(f_*\mathcal{O}_X)
\longrightarrow S
$$
of Constructions, Lemma \ref{constructions-lemma-canonical-morphism}.
\end{lemma}

\begin{proof}
The question is local on $S$, hence we may assume $S = \Spec(R)$
is affine. Let $h \in \Gamma(X, \mathcal{O}_X)$. We have to show that
$h$ satisfies a monic equation over $R$. Think of $h$ as a morphism
as in the following commutative diagram
$$
\xymatrix{
X \ar[rr]_h \ar[rd]_f & & \mathbf{A}^1_S \ar[ld] \\
& S &
}
$$
Let $Z \subset \mathbf{A}^1_S$ be the scheme theoretic image of $h$,
see Definition \ref{definition-scheme-theoretic-image}.
The morphism $h$ is quasi-compact as $f$ is quasi-compact and
$\mathbf{A}^1_S \to S$ is separated, see
Schemes, Lemma \ref{schemes-lemma-quasi-compact-permanence}.
By Lemma \ref{lemma-quasi-compact-scheme-theoretic-image} the
morphism $X \to Z$ is dominant. By
Lemma \ref{lemma-image-proper-scheme-closed} the morphism
$X \to Z$ is closed. Hence $h(X) = Z$ (set theoretically).
Thus we can use
Lemma \ref{lemma-image-proper-is-proper}
to conclude that $Z \to S$ is universally closed (and even proper).
Since $Z \subset \mathbf{A}^1_S$, we see that $Z \to S$ is affine
and proper, hence integral by Lemma \ref{lemma-integral-universally-closed}.
Writing $\mathbf{A}^1_S = \Spec(R[T])$ we conclude that
the ideal $I \subset R[T]$ of $Z$ contains a monic polynomial
$P(T) \in R[T]$. Hence $P(h) = 0$ and we win.
\end{proof}

\begin{lemma}
\label{lemma-normalization-in-integral}
Let $f : Y \to X$ be an integral morphism.
Then the normalization of $X$ in $Y$ is equal to $Y$.
\end{lemma}

\begin{proof}
By Lemma \ref{lemma-integral-universally-closed} this is a special case of
Lemma \ref{lemma-normalization-in-universally-closed}.
\end{proof}

\begin{lemma}
\label{lemma-normal-normalization}
Let $f : Y \to X$ be a quasi-compact and quasi-separated morphism
of schemes. Let $X'$ be the normalization of $X$ in $Y$. Assume
\begin{enumerate}
\item $Y$ is a normal scheme,
\item quasi-compact opens of $Y$ have finitely many irreducible components.
\end{enumerate}
Then $X'$ is a disjoint union of integral normal schemes. Moreover,
the morphism $Y \to X'$ is dominant and induces a bijection of
irreducible components.
\end{lemma}

\begin{proof}
Let $U \subset X$ be an affine open. Consider the
inverse image $U'$ of $U$ in $X'$.
Set $V = f^{-1}(U)$. By Lemma \ref{lemma-normalization-localization}
we $V \to U' \to U$ is the normalization of $U$ in $V$. Say
$U = \Spec(A)$. Then $V$ is quasi-compact, and hence has a finite number of
irreducible components by assumption. Hence
$V = \coprod_{i = 1, \ldots n} V_i$ is a finite disjoint union of
normal integral schemes by
Properties, Lemma \ref{properties-lemma-normal-locally-finite-nr-irreducibles}.
By Lemma \ref{lemma-normalization-in-disjoint-union}
we see that $U' = \coprod_{i = 1, \ldots, n} U_i'$,
where $U'_i$ is the normalization of $U$ in $V_i$.
By Properties, Lemma \ref{properties-lemma-normal-integral-sections}
we see that $B_i = \Gamma(V_i, \mathcal{O}_{V_i})$ is a normal domain.
Note that $U_i' = \Spec(A_i')$, where $A_i' \subset B_i$
is the integral closure of $A$ in $B_i$, see
Lemma \ref{lemma-integral-closure}. By
Algebra, Lemma \ref{algebra-lemma-integral-closure-in-normal}
we see that $A_i' \subset B_i$ is a normal domain.
Hence $U' = \coprod U_i'$ is a finite union of normal integral schemes
and hence is normal.

\medskip\noindent
As $X'$ has an open covering by the schemes $U'$ we conclude from
Properties, Lemma \ref{properties-lemma-locally-normal} that $X'$ is normal.
On the other hand, each $U'$ is a finite disjoint union of irreducible
schemes, hence every quasi-compact open of $X'$ has finitely many irreducible
components (by a topological argument which we omit). Thus $X'$
is a disjoint union of normal integral schemes by
Properties, Lemma \ref{properties-lemma-normal-locally-finite-nr-irreducibles}.
It is clear from the description of $X'$ above that $Y \to X'$
is dominant and induces a bijection on irreducible components
$V \to U'$ for every affine open $U \subset X$. The bijection of
irreducible components for the morphism $Y \to X'$
follows from this by a topological argument (omitted).
\end{proof}

\begin{lemma}
\label{lemma-nagata-normalization-finite-general}
Let $f : X \to S$ be a morphism. Assume that
\begin{enumerate}
\item $S$ is a Nagata scheme,
\item $f$ is quasi-compact and quasi-separated,
\item quasi-compact opens of $X$ have finitely many irreducible components,
\item if $x \in X$ is a generic point of an irreducible component,
then the field extension $\kappa(x)/\kappa(f(x))$ is finitely
generated, and
\item $X$ is reduced.
\end{enumerate}
Then the normalization $\nu : S' \to S$ of $S$ in $X$ is finite.
\end{lemma}

\begin{proof}
There is an immediate reduction to the case $S = \Spec(R)$
where $R$ is a Nagata ring by assumption (1). We have to show that
the integral closure $A$ of $R$ in $\Gamma(X, \mathcal{O}_X)$ is
finite over $R$. Since $f$ is quasi-compact by assumption (2) we can write
$X = \bigcup_{i = 1, \ldots, n} U_i$ with each $U_i$ affine.
Say $U_i = \Spec(B_i)$. Each $B_i$ is reduced by assumption (5)
and has finitely many minimal primes
$\mathfrak q_{i1}, \ldots, \mathfrak q_{im_i}$
by assumption (3) and
Algebra, Lemma \ref{algebra-lemma-irreducible}.
We have
$$
\Gamma(X, \mathcal{O}_X) \subset B_1 \times \ldots \times B_n
\subset
\prod\nolimits_{i = 1, \ldots, n}
\prod\nolimits_{j = 1, \ldots, m_i} (B_i)_{\mathfrak q_{ij}}
$$
the second inclusion by
Algebra, Lemma \ref{algebra-lemma-reduced-ring-sub-product-fields}.
We have $\kappa(\mathfrak q_{ij}) = (B_i)_{\mathfrak q_{ij}}$ by
Algebra, Lemma \ref{algebra-lemma-minimal-prime-reduced-ring}.
Hence the integral closure $A$ of $R$ in $\Gamma(X, \mathcal{O}_X)$
is contained in the product of the integral closures $A_{ij}$ of
$R$ in $\kappa(\mathfrak q_{ij})$. Since $R$ is Noetherian
it suffices to show that $A_{ij}$ is a finite $R$-module for each $i, j$.
Let $\mathfrak p_{ij} \subset R$ be the image of $\mathfrak q_{ij}$.
As $\kappa(\mathfrak q_{ij})/\kappa(\mathfrak p_{ij})$
is a finitely generated field extension by assumption (4),
we see that $R \to \kappa(\mathfrak q_{ij})$ is essentially of finite type.
Thus $R \to A_{ij}$ is finite by Algebra, Lemma
\ref{algebra-lemma-nagata-in-reduced-finite-type-finite-integral-closure}.
\end{proof}

\begin{lemma}
\label{lemma-nagata-normalization-finite}
Let $f : X \to S$ be a morphism. Assume that
\begin{enumerate}
\item $S$ is a Nagata scheme,
\item $f$ is of finite type,
\item $X$ is reduced.
\end{enumerate}
Then the normalization $\nu : S' \to S$ of $S$ in $X$ is finite.
\end{lemma}

\begin{proof}
This is a special case of
Lemma \ref{lemma-nagata-normalization-finite-general}.
Namely, (2) holds as the finite type morphism $f$ is quasi-compact
by definition and quasi-separated by
Lemma \ref{lemma-finite-type-Noetherian-quasi-separated}.
Condition (3) holds because $X$ is locally Noetherian by
Lemma \ref{lemma-finite-type-noetherian}. Finally, condition (4)
holds because a finite type morphism induces finitely generated
residue field extensions.
\end{proof}

\begin{lemma}
\label{lemma-relative-normalization-normal-codim-1}
Let $f : Y \to X$ be a finite type morphism of schemes with $Y$ reduced
and $X$ Nagata. Let $X'$ be the normalization of $X$ in $Y$.
Let $x' \in X'$ be a point such that
\begin{enumerate}
\item $\dim(\mathcal{O}_{X', x'}) = 1$, and
\item the fibre of $Y \to X'$ over $x'$ is empty.
\end{enumerate}
Then $\mathcal{O}_{X', x'}$ is a discrete valuation ring.
\end{lemma}

\begin{proof}
We can replace $X$ by an affine neighbourhood of the image of $x'$.
Hence we may assume $X = \Spec(A)$ with $A$ Nagata.
By Lemma \ref{lemma-nagata-normalization-finite}
the morphism $X' \to X$ is finite. Hence we can write
$X' = \Spec(A')$ for a finite $A$-algebra $A'$.
By Lemma \ref{lemma-normalization-is-normalization}
after replacing $X$ by $X'$
we reduce to the case described in the next paragraph.

\medskip\noindent
The case $X = X' = \Spec(A)$ with $A$ Noetherian.
Let $\mathfrak p \subset A$ be the prime ideal
corresponding to our point $x'$. Choose $g \in \mathfrak p$
not contained in any minimal prime of $A$ (use prime avoidance
and the fact that $A$ has finitely many minimal primes, see
Algebra, Lemmas \ref{algebra-lemma-silly} and
\ref{algebra-lemma-Noetherian-irreducible-components}).
Set $Z = f^{-1}V(g) \subset Y$; it is a closed subscheme of $Y$.
Then $f(Z)$ does not contain any generic point by choice of $g$
and does not contain $x'$ because $x'$ is not in the image of $f$.
The closure of $f(Z)$ is the set of specializations of points
of $f(Z)$ by Lemma \ref{lemma-reach-points-scheme-theoretic-image}.
Thus the closure of $f(Z)$ does not contain $x'$ because
the condition $\dim(\mathcal{O}_{X', x'}) = 1$ implies only
the generic points of $X = X'$ specialize to $x'$.
In other words, after replacing $X$ by an affine open
neighbourhood of $x'$ we may assume that $f^{-1}V(g) = \emptyset$.
Thus $g$ maps to an invertible global function on $Y$ and
we obtain a factorization
$$
A \to A_g \to \Gamma(Y, \mathcal{O}_Y)
$$
Since $X = X'$ this implies that $A$ is equal to the integral
closure of $A$ in $A_g$. By
Algebra, Lemma \ref{algebra-lemma-integral-closure-localize}
we conclude that $A_\mathfrak p$ is the integral closure
of $A_\mathfrak p$ in $A_\mathfrak p[1/g]$.
By our choice of $g$, since $\dim(A_\mathfrak p) = 1$
and since $A$ is reduced we see that $A_\mathfrak p[1/g]$
is a finite product of fields (the product of the residue fields
of the minimal primes contained in $\mathfrak p$). Hence $A_\mathfrak p$ is
normal (Algebra, Lemma \ref{algebra-lemma-characterize-reduced-ring-normal})
and the proof is complete. Some details omitted.
\end{proof}






